\FloatBarrier
\label{sec:DTData}

For this work, atmospheric muon data was recorded with the \ac{DT} chamber at the reactivated test stand in Aachen (sec.\ \ref{sec:CMS-AachenTestStand}). This section should present the data and the analysis strategy to reconstruct muon tracks without an external trigger. The triggerless reconstruction is essential, because the developed readout chain for the scintillation detector (sec.\ \ref{sec:Scint}) should be validated using the \ac{DT} chamber hits as reference, therefore using the scintillator as a trigger for the chamber would imply a bias.
\newline The reconstruction strategy is guided by the local \ac{CMS} \ac{DT} reconstruction \cite{DT-LocalReco} and calibration \cite{DT-OfflineCalibration}, as well as other work within the \ac{DT} group concerning \ac{DT} hit reconstruction \cite{MSc-Alcalde,CMS-NOTE-2008-017}. The reconstruction was implemented into a custom standalone code framework \cite{dt-scint-analysis-tools-Main}, which allows the analysis of the \ac{DT} chamber data as well as the scintillator data (sec.\ \ref{sec:ScintData}). Integration with the central \ac{CMS} \ac{DT} reconstruction framework might be desired in the future, to exploit the evolved algorithms. However, the aim of this thesis is a proof-of-concept measurement with the \ac{DT} chamber and the scintillator. Therefore, a less evolved reconstruction strategy inside this self-developed framework is considered sufficient.





\FloatBarrier
\section{Timing calibration with testpulses}
%\begin{itemize}
%    \item Importance of timing calibration in DT
%    \item Refer to CMS DT calibration procedure
%    \item Show testpulse response
%    \item Discuss different testpulse cable length in theta SL
%    \item Show derived timing corrections for individual channels
%\end{itemize}
\label{sec:DTData-Testpulse}
Precise time measurement of the drift cell hits is essential for a successful \ac{DT} reconstruction, as the measured timestamps are directly used to calculate the hit position within the drift volume (sec.\ \ref{sec:GaseousDetectors-DriftChamber}). Therefore, it must be ensured that the readout of all channels of the \ac{DT} chamber is actually synchronized \cite{DT-OfflineCalibration}. 

\FloatBarrier
\subsection{Testpulse generation}
\label{sec:DTData-Testpulse-Generation}
\FloatBarrier
\paragraph{Synchronous testpulse generation}
As described, the minicrate can generate synchronous testpulses for the chamber's \acp{FEB}, which sends back a response to the \acp{OBDT} (sec.\ \ref{sec:CMS-DTMinicrate}). The issuing of testpulses is controlled by the \ac{DT} backend. It can send a command to all \acp{OBDT} to generate a testpulse at a given $BX$ time in an orbit (a new orbit starts at $BX=0$ and $TDC=0$, as discussed in sec.\ \ref{sec:CoordinateFrame-TimestampFormat}). Therefore, the time when the testpulse is issued is precisely known.
\FloatBarrier
\paragraph{Testpulse latency correction}
The testpulse generators of the \ac{OBDT} phi and theta boards have different analog circuitry which produce the actual testpulses (sec.\ \ref{sec:CMS-Phase2Readout}). This results in a latency difference between the boards, which is however constant, as it originates from the (static) circuit. In order to eliminate this timing offset between \ac{OBDT} flavors, a correction should be applied to the recorded testpulse hit timestamps. A measurement of the time difference of the generated testpulse output with an oscilloscope yields
\begin{equation}
    T_\text{testpulse, theta} - T_\text{testpulse, phi} = (116 \pm 1)\;\si{\nano\second}.
    \label{eq:DTData-TestpulseLatency}
\end{equation}
A screenshot of the oscilloscope measurement can be found in the appendix in fig.\ \ref{fig:DTData-Testpulse-LatencyDifference}. This latency correction value is applied to the recorded timestamps of the testpulse hits.
\newline Additionally, the testpulse cable lengths from the minicrate to the \acp{FEB} in the chamber are not all the same. To account for latency differences due to the different lengths, the difference in signal propagation time was measured with an oscilloscope and applied to the affected channels as well.

\FloatBarrier
\subsection{Testpulse response and timing calibration}
\label{sec:DTData-Testpulse-Calibration}
A dataset with activated testpulses is recorded, where all \ac{DT} chamber channels are unmasked in the \ac{OBDT} to record the testpulse responses. Since the time of the testpulse generation is fixed w.r.t.\ the orbit counter, one can use the time relative to the beginning of the respective orbit for this analysis, which is defined as
\begin{equation}
    T_\text{orbit} = T(OC,\;BX,\;TDC) - T(OC,\;BX=0,\;TDC=0).
\end{equation}
For the presented testpulse run, the pulses were commanded to be sent at $BX=50$ of each orbit, translating to $T_\text{orbit}=\SI{1600}{\tu}$. The testpulse response should come later than this time, as one has to expect several delay sources for the received response: The latency of the testpulse generation circuitry, the propagation time via the testpulse cables, the \ac{FEB} intrinsic response time as well as the frontend cable lengths of the \acp{FEB} going back to the \ac{OBDT}.
\FloatBarrier
\paragraph{Testpulse response}
The timestamps of all received testpulses for the each channel are averaged to obtain one testpulse time value per channel $\left\langle T_\text{orbit}(\text{SL, Ly, Wi}) \right\rangle$. Fig.\ \ref{fig:DTData-TestpulseResponse} shows the testpulse response for all channels of the \ac{DT} chamber. The timing uncertainties are calculated from the digitization step of the \ac{OBDT} (\SI{0.78}{\nano\second}) and the uncertainty of the latency correction discussed above (sec.\ \ref{sec:DTData-Testpulse-Generation}).
\begin{figure}[htb]
    \centering
    \includegraphics[width=\textwidth]{dt_tp/t_orbit.pdf}
    \caption{Mean timestamp $\left\langle T_\text{orbit}(\text{SL, Ly, Wi}) \right\rangle$ of the testpulse responses for all individual channels.}
    \label{fig:DTData-TestpulseResponse}
\end{figure}
\newline One can observe that all read-out channels do actually receive a testpulse response. From this it is concluded that there are no dead channels on the \ac{OBDT} side and the installed minicrate is working correctly.
\newline As expected, the testpulse response is recorded at timestamps greater than the time $T_\text{orbit}=\SI{1600}{\tu}$ when the testpulse generation is commanded. The recorded response timestamps lie in the range $T_\text{orbit}\approx 1705-1720 \;\si{\tu}$. While for the phi \acp{SL} 1 and 3 the interval of all response timestamps is comparable, the response of all channels of the theta \ac{SL} 2 is systematically later by $\approx 10\;\si{\tu}$ ($\approx 7 \;\si{\nano\second}$). One possible reason is the length of the \ac{FEC} cables, which contributes to the response timing. Since the minicrate is mounted on the chamber side with the \ac{SL} phi \acp{FEC}, the phi cables are significantly shorter ($\approx 0.3-1.6\;\si{\meter}$) than the theta cables ($\approx 3\;\si{\meter}$). A cable length difference of $\approx\SI{2}{\meter}$ could in fact explain a timing difference in the order of magnitude that is observed.\footnote{With the rough assumption of an electrical signal propagation velocity of $\approx\SI{30}{\centi\meter/\nano\second}$, one obtains an expected delay of \SI{6}{\nano\second} for a cable length difference of \SI{2}{\meter}.} Since this difference in cable length is in the signal path, it is right to correct for this in the timing calibration.
\newline Another observation is the difference in the channel-to-channel variations between the phi and theta \acp{SL}. While the testpulse responses for \ac{SL} 2 (theta) are contained within a window of only \SI{5}{\tu} (\SI{4}{\nano\second}), the channels vary by up to \SI{10}{\tu} (\SI{8}{\nano\second}) for the phi \acp{SL} 1 and 3. The \ac{FEC} cables of the phi \acp{SL} are not uniform in length, as it is the case for the theta (see mentioned values above). This could again explain the order of magnitude in fluctuation that is observed. One would expect that the cable length affects groups of channels that are connected to the same cable in the same way. However, one can also observe significant variations between channels of the same wire in different layers, which must have a different cause.
\FloatBarrier
\paragraph{Extracted timing correction}
The relative timing correction for all channels can now be calculated. For the correction, it is decided to use the global average of testpulse times $\left\langle T_\text{orbit} \right\rangle$ as a reference time. One can estimate a relative timing correction $T_\text{corr}$ for the individual channels by subtracting the channel timestamp by this reference time:
\begin{equation}
    T_\text{corr}(\text{SL, Ly, Wi}) = \left\langle T_\text{orbit}(\text{SL, Ly, Wi}) \right\rangle - \left\langle T_\text{orbit} \right\rangle.
\end{equation}
Note that because the latency correction was performed (eq.\ \ref{eq:DTData-TestpulseLatency}), the obtained timing correction covers the full readout path (from \ac{FEB} over the cables to the \acp{OBDT} and the path to the \ac{DT} backend), but not the variations in the testpulse generation. The extracted corrections are shown in fig.\ \ref{fig:DTData-TestpulseCorrections}.
\begin{figure}[htb]
    \centering
    \includegraphics[width=\textwidth]{dt_tp/t_corr.pdf}
    \caption{Calculated relative timing corrections $T_\text{corr}(\text{SL, Ly, Wi})$ for all individual channels.}
    \label{fig:DTData-TestpulseCorrections}
\end{figure}
\newline The obtained timing corrections lie in the range of $-6\;\si{\nano\second}$ to $+7\;\si{\nano\second}$. In comparison, the \ac{CMS} \ac{DT} offline calibration note \cite{DT-OfflineCalibration} states usual timing calibration values of up to $\pm 9\;\si{\nano\second}$. Thererfore, the result from the testpulse run presented above seems to yield calibration values in the expected range.





\FloatBarrier
\section{Drift cell hits}
%\begin{itemize}
%    \item Show hit occupancies of all wires
%    \item Give list of masked (noisy) and dead wires
%    \item Discuss rate
%    \item Highlight and discuss feature of secondary hits in range 0-400 ns
%    \item Apply dead time for individual channels
%\end{itemize}
\label{sec:DTData-CellHits}
In the following, a dataset of atmospheric muons recorded with the \ac{DT} chamber at the test stand is presented. Tab.\ \ref{tab:DTData-DatasetInfo} presents some basic facts about the recorded dataset. Before the reconstruction algorithm is described and its results are showcased, the raw data (i.e.\ the drift cell hit timestamps) should be presented and discussed in this section.
\begin{table}[htb]
    \centering
    \begin{tabular}{|c|c|c|} 
        \hline % size = 13.6 gb * 0.47 (approx. dt fraction on dumpfile) = 
        Measurement duration & Total \ac{DT} cell hits & Raw data size \\ \hline
        $\SI{19599}{\second}\approx\SI{5}{\hour}\;\SI{27}{\min}$ & \num{297312727} & $\SI{6.4}{\giga\byte}$ \\ \hline
    \end{tabular}
    \caption{Basic facts about the presented \ac{DT} chamber dataset with atmospheric muons.}
    \label{tab:DTData-DatasetInfo}
\end{table}

\FloatBarrier
\subsection{Cell occupancy}
\label{sec:DTData-CellHits-Occupancy}
As the \acp{OBDT} in the minicrate stream out all hits of the chamber and do not apply any preselection (sec.\ \ref{sec:CMS-DTMinicrate}), the recorded \ac{DT} hits are expected to include both real hits caused by ionization events in the cell as well as noise, which is discussed in the following sections. Note that the data presented here has the timing correction from above (sec.\ \ref{sec:DTData-Testpulse-Calibration}) already applied.
\FloatBarrier
\paragraph{Measured cell occupancies}
The full dataset contains a total of \num{297312727} drift cell hits. Fig.\ \ref{fig:DTData-ChamberRates} shows the cell occupancies of all individual wires of the chamber. The occupancy is observed to be relatively uniform for all wires and is $\mathcal{O}(10^5)$ hits per cell for the full recorded dataset. With the known recording time of the data, one can calculate a mean event rate for each cell. Tab.\ \ref{tab:DTData-ChamberRates} presents the cell rates for different groups of cells within the chamber. Over the full chamber, the average cell rate is measured as $(25.111 \pm 0.001)\;\si{\hertz}$.
\begin{figure}[htb]
    \centering %% VALUES: WITH FULL DATASET
    \includegraphics[width=\textwidth]{large_combined_data_plots/53_comb__DT_CORR_HITS_SPECIFIC_OCCUPANCY.pdf}
    \caption{Occupancies of all individual drift cells of the \ac{DT} chamber.}
    \label{fig:DTData-ChamberRates}
\end{figure}
\begin{table}[htb]
    \centering
    \begin{tabular}{|c|c|}  %% VALUES: WITH FULL DATASET
        \hline
        Chamber subset & Average cell rate \\ \hline
        \ac{SL} 1 (phi) & $(26.676\pm 0.003)\;\si{\hertz}$ \\
        \ac{SL} 2 (theta) & $(22.940 \pm 0.002)\;\si{\hertz}$ \\
        \ac{SL} 3 (phi) & $(25.976 \pm 0.003)\;\si{\hertz}$ \\ \hline
        \acp{SL} 1, 3 (phi) & $(26.329 \pm 0.002)\;\si{\hertz}$ \\ \hline
        Full chamber & $(25.111 \pm 0.001)\;\si{\hertz}$ \\ \hline
    \end{tabular}
    \caption{Averaged \ac{DT} cell rates for different subsets of the chamber. The dead cells (tab. \ref{tab:DTData-DeadWires}) were not considered for the rate averaging. Statistical uncertainties are given for each value.}
    \label{tab:DTData-ChamberRates}
\end{table}
\FloatBarrier
\paragraph{Noisy and dead cells}
In fig.\ \ref{fig:DTData-ChamberRates} few cells with low occupancy are visible. In principle, there are two different kinds of differently behaving cells. First, there are cells which do not produce a signal in the \ac{FEB}. These cells are called dead cells. They might originate from either broken electrical or high-voltage connections inside the chamber or problems on the \ac{FEB} itself. On the other hand, there might be noisy cells which exhibit much higher hit rates than one would expect. In the case of the \ac{DT} chamber at the test stand, no noisy cells were found. Tab.\ \ref{tab:DTData-DeadWires} gives a comprehensive list of the observed dead wires. Note that although the layer 1 of each \ac{SL} has one additional cell compared to the other layers (tab.\ \ref{tab:MB1Chamber-Cells}), this additional cell is not connected to the readout. Therefore, the observed zero occupancy in the right most cell of layer 1 of each \ac{SL} is expected. In total, 13 dead wires plus 3 unconnected wires are observed in the \ac{DT} chamber at the test stand. The most dead wires are found in \ac{SL} 2, while \ac{SL} 1 does not have any dead wires. Additionally, one can observe in fig.\ \ref{fig:DTData-ChamberRates} that the wire numbering of layer 3 in each \ac{SL} starts at 1 and not at 0 (cf.\ tab.\ \ref{tab:MB1Chamber-Cells}).
\begin{table}[htb]
    \centering
    \begin{tabular}{|c|c|c|c|}  %% VALUES: WITH FULL DATASET
        \hline
        \ac{SL} & Ly & Wi & Type \\ \hline
        % ro_ch=8 masked chs: 115, 119, 121, 123 (fec 7A, probably higher than use chs)
        1 (phi) & 1 & 49 & Not connected \\ %\hline
        % ro_ch=14 masked chs: none
        2 (theta) & 0 & 2 & Dead \\ 
        2 (theta) & 0 & 5 & Dead \\ 
        2 (theta) & 0 & 13 & Dead \\ 
        2 (theta) & 0 & 35 & Dead \\ 
        2 (theta) & 0 & 50 & Dead \\ 
        2 (theta) & 0 & 51 & Dead \\ 
        2 (theta) & 1 & 4 & Dead \\ 
        2 (theta) & 1 & 5 & Dead \\ 
        2 (theta) & 1 & 44 & Dead \\ 
        2 (theta) & 1 & 45 & Dead \\ 
        2 (theta) & 1 & 57 & Not connected \\ %\hline
        % ro_ch=10 masked chs: 124 (fec 7A, probably higher than use chs)
        3 (phi) & 0 & 10 & Dead \\ 
        3 (phi) & 1 & 41 & Dead \\ 
        3 (phi) & 1 & 49 & Not connected \\ 
        3 (phi) & 2 & 26 & Dead \\ \hline
    \end{tabular}
    \caption{List of dead and unconnected wires of the \ac{DT} chamber at the test stand.}
    \label{tab:DTData-DeadWires}
\end{table}
\FloatBarrier
\paragraph{Comparison with CMS}
Different \ac{CMS} publications present hit rates which lie in the range $\approx 30-60\;\si{\hertz}$ per cell \cite{CMS-Main,CMS-DT-TestBeam-2004}. These rates are slightly higher than what is observed in the recorded dataset from the test stand. The presumable higher noise in the \ac{CMS} data may be related to older readout electronics that were used at the time, as the new minicrate was not available yet. A lower selected \ac{FEB} threshold voltage might also lead to an increased hit rate. The threshold of the cited \ac{CMS} values is not known, but might have been lower than the selected \SI{35}{\milli\volt} (tab.\ \ref{tab:MB1Chamber-Settings}) used in this work. Although the hit rates in \ac{CMS} publications are larger, the order of magnitude $\mathcal{O}(10^1\;\si{\hertz})$ is still the same as what is observed in the measurements of this thesis.
\FloatBarrier  %% CORRECT SIM EXP VALUES
\paragraph{Comparison with cosmic rate} A comparison of the measured cell rates to the rate expectation solely from atmospheric muons should be made in this paragraph. In principle, the expected rate can be calculated, if one solves the solid angle and area integrals of the muon flux for particles above the energy detection threshold (eq.\ \ref{eq:CosmicMuon-DetectorRate}), as was discussed in sec.\ \ref{sec:AtmosphericMuons-Cosmics}. However, for non-spherical geometries this integration is non-trivial.
\newline Alternatively, it is possible to estimate the rate with a \ac{MC} simulation. The developed analysis code includes this possibility of simulation: It can randomly generate muon tracks which follow the atmospheric muon rate and angular distribution and propagate them through the detector geometry. Note that the simulation does not model the physics of the muon interacting with the detector material, but simply simulates the geometric tracks and detects intersections with sensitive cell volume. However, the simulation respects the finite hit efficiency of the \ac{DT} cells (sec.\ \ref{sec:CMS-DTPerformance}). Details the assumptions and simplifications of this simulation are discussed in appendix \ref{sec:Simulation}.
\newline The simulation yields an average atmospheric muon rate of $(14.342 \pm 0.006)\;\si{\hertz}$ for a phi drift cell and $(12.302 \pm 0.007)\;\si{\hertz}$ for a theta drift cell (tab.\ \ref{tab:Simulation-ChamberRates}). These values are approximately $50\%$ smaller than the observed cell rates of the \ac{DT} chamber at the test stand. The higher observed cell rate could be explained with the presence of noise hits in the real detector, as the simulation performed here does not include any modeling of noise. Even though there is a deviation, the order of magnitude of $\mathcal{O}(10^1\;\si{\hertz})$ per cell is still similar to what is observed in data. Together with the fact that the cell rates cited in \ac{CMS} publications for cosmics are also in a similar order of magnitude, it is concluded that the recorded \ac{DT} data seems plausibe and that the chamber seems to operate nominally.
\FloatBarrier
\paragraph{Time uniformity}
Fig.\ \ref{fig:DTData-Hist-T} presents a histogram of all drift cell hit timestamps in the recorded dataset. As one would expect from cosmics and noise, the distribution is uniform in time. At the edges the histogram count is lower, however is purely a binning effect. Would there have been spikes in the timestamp distribution, this would have hinted some problem with the readout system, as uniform illumination from atmospheric muons and a constant noise rate are expected.
%\begin{figure}[htb]
%    \centering
%    \includegraphics[width=0.7\textwidth]{large_combined_data_plots/53_comb__DT_CORR_HITS_HIST_ts.pdf}
%    \caption{Histogram of the recorded drift cell hit timestamps $T$.}
%    \label{fig:DTData-Hist-T}
%\end{figure}
\begin{SCfigure}[50][htb]
    \centering
    \includegraphics[width=0.49\textwidth]{large_combined_data_plots/53_comb__DT_CORR_HITS_HIST_ts.pdf}
    \caption{Histogram of the recorded \ac{DT} hit timestamps $T$.}
    \label{fig:DTData-Hist-T}
\end{SCfigure}

%As discussed in sec.\ \ref{sec:AtmosphericMuons-Cosmics}, the detector rate can be calculated by integration of the muon flux (cf.\ eq.\ \ref{eq:CosmicMuon-DetectorRate}). The solution of the solid angle and area integral is non-trivial, therefore 
%A rough estimation for the vertical cosmic hit rate onto one drift cell yields $\approx\SI{10}{\hertz}$ ($\approx\SI{8.5}{\hertz}$) for the phi (theta) \acp{SL}. However, cosmic particles at higher angles are not considered in this estimation as well as the influence of noise.

\FloatBarrier
\subsection{Secondary hits}
\label{sec:DTData-CellHits-SecondaryHits}
If one considers all hits of one single drift cell, sorts them by timestamp and plots the time difference $\Delta T_\text{cell}$ between these hits, one can observe some characteristic timing features of the \ac{DT} cell hits. Fig.\ \ref{fig:DTData-CellHitDifference} shows the distribution for this time difference $\Delta T_\text{cell}$ above, cumulated for all cells. Note that the $x$-axis limit extends far beyond the presented range, but the discussed features are already visible in this limited range.
\begin{figure}[htb]
    \centering
    \includegraphics[width=0.7\textwidth]{dt_hits/Hit_difference.pdf}
    \caption{Histogram of the time difference $\Delta T_\text{cell}$ between hits of a single drift cell, cumulatively plotted for all cells.}
    \label{fig:DTData-CellHitDifference}
\end{figure}
\FloatBarrier
\paragraph{Poisson background}
One can observe that the time difference between hits decreases exponentially. This behavior continues to the right beyond the plotted range. Both cosmic muon hits and noise are random events which can be described by a Poisson distribution. From statistics it is known that the time difference between Poisson events is exponentially distributed \cite{Book-KolanoskiWermes}. Therefore, this behavior is expected and confirms the random nature of the recorded cell hits.
\FloatBarrier
\paragraph{Secondary hits}
In the region $\approx 0-500\;\si{\nano\second}$, another feature is visible which dominates over the exponentially decreasing behavior. Apparently, there are hits which are detected in this short time interval after the initial hit. The observed structure of these secondary hits should be analyzed in the following. For this, the Poissonian background due to random noise events should be subtracted, to analyze the shape of secondary hits in time. For this, an exponential function is fitted to the exponentially decreasing part of the histogram (fit range: $\Delta T_\text{cell} \geq \SI{1000}{\nano\second}$). Fig.\ \ref{fig:DTData-CellHitDifference-ExpFit} in the appendix shows the fit result.
\newline Fig.\ \ref{fig:DTData-CellHitDifference-NoBg} shows the resulting distribution after extrapolating the exponential "background" to the left and subtracting it. A structure of two peaks can be observed, where the second peak is relatively narrow, while the first peak is wider. Also note that one can observe a dead time of $\approx\SI{40}{\nano\second}$ in the plot, which might originate from the \acp{FEB} in the chamber or the \acp{OBDT}.
\newline Note that this structure of secondary peaks was not discovered in this work, but is a known effect. Two \ac{CMS} Notes \cite{CMS-NOTE-2008-003,CMS-NOTE-2010-004} also show a distribution of secondary hits which looks similar in shape. In fig.\ \ref{fig:DTData-CellHitDifference-CMSNote} the observed distribution of the secondary hits from one of these publications is shown for comparison. The distribution matches well the double-peak structure which is observed in the measurement from this thesis with the background subtracted.
\begin{figure}[htb]
    \centering
    \hfill
    \begin{subfigure}[c]{0.49\textwidth}
        \centering
        \includegraphics[width=\textwidth]{dt_hits/Hit_difference_nobg.pdf}
        \caption{Time distribution of secondary hits after extrapolating and extracting the exponentially decreasing noise background.}
        \label{fig:DTData-CellHitDifference-NoBg}
    \end{subfigure}
    \hfill
    \begin{subfigure}[c]{0.49\textwidth}
        \centering
        \includegraphics[width=\textwidth]{Hit_difference_CMSNote.png}
        \caption{Time distribution of secondary hits measured by \ac{CMS} \cite{CMS-NOTE-2008-003}.}
        \label{fig:DTData-CellHitDifference-CMSNote}
    \end{subfigure}
    \hfill
    \caption{Time distribution of secondary hits in this work compared to a \ac{CMS} measurement.}
    \label{fig:DTData-CellHitDifference-SecondaryHits}
\end{figure}
\FloatBarrier
\paragraph{Physical origin}
The observation that the secondary peaks lie in the time range of $\approx 0-500\;\si{\nano\second}$ already hints that these hits might have a physical cause, and are not created by reflections in the readout electronics, as the electron drift time in a \ac{DT} cell should lie in this order of magnitude (cf.\ eq.\ \ref{eq:DT-MaxDriftTime}). Two physical effects play a role in the creation of secondary peaks: When the electrons in the drift cell arrive at the wire, an avalanche forms and the charge is amplified and read out. In this avalanche, photons may be created, whose movement is not constrained by the electric field in the cell. Therefore, the photons can hit the chamber walls and create a secondary ionization event through the photoelectric effect. The created electrons from this secondary ionization drift again towards the wire, creating a secondary hit after the first one. Since the drift cell is constrained in size, the timestamps of the secondary hits also have a maximum. The \ac{DT} cell includes two electrodes apart from the anode wire (cf.\ fig.\ \ref{fig:DT-Cell}): The field shaping electrodes which are relatively close to the wire, and the cathode strips which are roughly \SI{21}{\milli\meter} (i.e.\ half of the cell width) away from the wire. When secondary ionization happens on one of these electrodes, the photo-electrons can reach the wire and generate a signal. When the ionization takes place somewhere else on the cell wall, the electrons are most likely lost and can not reach the wire, due to the shape of the electric field lines. The two-peak structure observed in the data can be explained with photo-electrons originating from the field shaping electrodes (first wide peak) and the cathode strips (second narrow peak). These peaks are in the following referred to as photo-peaks \cite{CMS-NOTE-2008-003}.
\newline The second physical effect responsible for secondary hits are so-called $\delta$-electrons. The energy loss of the passing muon is Landau-distributed, allowing larger energy losses in few collisions, which lead to high energy electrons in the gas volume (sec.\ \ref{sec:Interactions-Ionization}). These electrons may ionize the gas themselves, which can also contribute to secondary hits at the wire. Depending on the movement direction of the $\delta$-electron, not only after-pulses, but also pre-pulses may be created. The effect of $\delta$-electrons to the double peak structure is considered small compared to the photo-peaks, according to the \ac{CMS} Note \cite{CMS-NOTE-2008-003}.

\FloatBarrier
\subsection{Drift velocity estimation with photo-peak}
\label{sec:DTData-CellHits-DriftVelocityEstimation}
The second photo-peak can be used to estimate the drift velocity of the electrons in the \ac{DT} cell, without the need of an external trigger or a track reconstruction. This method was introduced in \cite{CMS-NOTE-2008-003}, where it was however only used for relative drift velocity measurements, not the absolute value.
\FloatBarrier
\paragraph{Photo-peak method}
As discussed above, the second (narrower) photo-peak originates from photo-electrons created at the cathode strip, which is located on the cell walls at roughly $d/2=\SI{21}{\milli\meter}$ distance from the wire. Assuming that the avalanche takes place very close to the wire and neglecting the photon travelling time, one can consider the peak position of the photo-peak as the drift time of the created photo-electron from the cathode strip to the wire. A parabolic fit to the second photo-peak is presented in fig.\ \ref{fig:DTData-CellHitDifference-PeakFit}. From the goodness-of-fit value and the residual plot it is concluded that the peak fit is acceptable, even though it is clear that the peak is asymmetric, and the parabolic fit therefore only is a first-order estimation.
\begin{figure}[htb]
    \centering
    \includegraphics[width=0.7\textwidth]{dt_hits/Hit_difference_nobg_peakfit.pdf}
    \caption{Parabolic fit to the second narrow peak in the time difference distribution of secondary hits.}
    \label{fig:DTData-CellHitDifference-PeakFit}
\end{figure}
\newline With the obtained photo-peak position $t_\text{pp}=b=(412.12 \pm 0.24)\;\si{\nano\second}$ and an uncertainty on the drift distance of $\sigma_{d/2}\approx\SI{100}{\micro\meter}$ (which was the requirement during the chamber construction \cite{CMS-Main}) one can estimate the drift velocity as
\begin{equation}
    v_D=\frac{d/2}{t_\text{pp}}\approx ( 51.0 \pm 0.2 )\;\si{\frac{\micro\meter}{\nano\second}}.
    \label{eq:PhotoPeak-DriftVelocityEstimation}
\end{equation}
Of course this value assumes a constant drift velocity over the full drift distance, which in reality is not the case. Another unregarded fact is the spatial extent of the cathode strip. This means that photo-electrons may also drift towards the wire from a different height, leading to a longer drift distance than the assumed half cell width $d/2$ and therefore a longer drift time. It therefore becomes clear that with this method one can only measure an effective drift velocity under the assumptions above, and that one needs to correct for the described geometric effect of longer drift distances.
\newline Within the scope of this thesis, no further simulation studies are performed, which would however be needed to account for this effect and to derive an average (effective) drift distance $x_{D,\text{eff}} > d/2$ which could average out this geometric effect and would allow the calculation of the drift velocity as $v_{D}=x_{D,\text{eff}} / t_\text{pp}$. 
\FloatBarrier
\paragraph{Comparison to CMS}
The obtained value for the drift velocity above (eq.\ \ref{eq:PhotoPeak-DriftVelocityEstimation}) is smaller than the expectation in the \ac{DT} system of $\approx\SI{54.5}{\micro\meter/\nano\second}$ (cf.\ eq.\ \ref{eq:DT-DriftVelocity}). However, the described geometric effects can explain the smaller value, although quantitative studies were not performed yet. Still one can consider the measurement of a drift velocity with this method a success, as the resulting velocity lies in the right order of magnitude of $\mathcal{O}(\SI{50}{\micro\meter/\nano\second})$.
\newline Note that in the CMS Note \cite{CMS-NOTE-2008-003} the photo-peak method was only used to estimate the relative change in drift velocity, never the absolute value. If one roughly estimates the photo-peak position from fig.\ \ref{fig:DTData-CellHitDifference-CMSNote} as $t_\text{pp}\approx\SI{510}{\tu}$ (\SI{398}{\nano\second}), one would obtain a drift velocity of $v_D\approx\SI{52.8}{\micro\meter/\nano\second}$ with the formula in eq.\ \ref{eq:PhotoPeak-DriftVelocityEstimation}. This value is also lower than the \ac{CMS} expectation (eq.\ \ref{eq:DT-DriftVelocity}).
\newline From this one can conclude that the photo-peak method is not capable of measuring the drift velocity in terms of absolute value without calibration. Still, the obtained drift velocities lie in the expected order of magnitude. As discussed, this absolute calibration may include simulation studies for the effective drift distance $x_{D,\text{eff}}$ of the photo-electrons. Studies at the \ac{DT} test stand involving a reference measurement with the \ac{VDC} monitoring hardware \cite{DT-VDC-System} could be imagined in the future.

\FloatBarrier
\subsection{Cell dead time}
\label{sec:DTData-CellHits-DeadTime}
As discussed above, there are occasions where secondary hits in the \ac{DT} cells can occur (sec.\ \ref{sec:DTData-CellHits-SecondaryHits}). If a track reconstruction would be attempted with these hits, one will obtain tracks which do not originate from a passing muon. As these secondary hits usually originate from photons escaping the electron amplification avalanche near the wire, the easiest way to suppress them is to apply a dead time to each individual \ac{DT} cell.  The dead time covers the time interval where the secondary hits occur. Of course this is not necessarily helpful in the case of $\delta$-electrons, as they also create pre-pulses and not only after-pulses. In the case of a pre-pulse, the dead time would lead to a loss of the "real" hit and therefore a wrong hit selection.
\newline Since this analysis does not aim at optimized precision but rather at a proof-of-concept measurement of the \ac{DT} chamber in combination with the scintillator readout, this small bias is accepted. In the following, a dead time of \SI{600}{\tu} (\SI{468}{\nano\second}) was applied to all individual cells to suppress the effect of secondary hits.




\FloatBarrier
\section{Superlayer hit clustering}
%\begin{itemize}
%    \item Triggerless reconstruction: Need to consider all recorded hits as potential muon signature
%    \item Cluster hits within maximum drift time window
%    \item Pattern rate discussion
%    \item Compare to rates of fake patterns (to estimate noise influence)
%\end{itemize}
\label{sec:DTData-SLHitPatterns}
The first step in the reconstruction of tracks in the \ac{DT} chamber is the selection of cell hit candidates, which may be part of a track. Since the reconstruction is triggerless, no external trigger signal can be used for this purpose. Instead, all recorded cell hits must be treated as potential track hit candidates and consecutively filtered out. Hit selection in this step is based on geometric shape of consecutive hits and their temporal distance.
\FloatBarrier
\subsection{Track candidate patterns}
\label{sec:DTData-SLHitPatterns-TrackPatterns}
\FloatBarrier
\paragraph{SL track patterns}
A muon track leaves ionization charges and consequently a hit in cells of the different layers of the \ac{DT} chamber. Depending on the track angle, different cell hit patterns can be produced. Fig.\ \ref{fig:DTData-SLPatterns} shows possible hit patterns in one \ac{SL}, under the assumption that four cells are triggered. It is clear that depending on the inclination of the muon track, different patterns can be produced. Note that in reality not always all cells are hit due to detector inefficiencies.
\begin{figure}[htb]
    \centering
    \includegraphics[width=0.9\textwidth]{SL-Patterns.pdf}
    \caption{Possible \ac{SL} hit patterns created by muon tracks. Only the patterns $0-5$ are considered for the analysis in this thesis.}
    \label{fig:DTData-SLPatterns}
\end{figure}
\FloatBarrier
\paragraph{Pattern timing}
Depending on the position of the muon track within each cell, the distance from the wire that the electrons need to drift along varies, which results in the hits of the pattern not having the same timestamps (i.e.\ time of collection of the electrons at the wire). However, one can formulate an upper bound of the timestamp differences of cell hits of a pattern, which might belong to a muon track. This upper bound is given by the maximum drift velocity, since the cell size sets limits the drift times of all cells:
\begin{equation}
    \Delta T_\text{max} = \max \left\{ \left | T_{ly=i} - T_{ly=j} \right | ; \; i,j=0,1,2,3 \right\} \leq t_\text{drift,max} \approx \SI{385}{\nano\second}.
    \label{eq:DTData-PatternTimeWindow}
\end{equation}
Here, $T_{\text{Ly}\; i}$ are the hit timestamps of the respective cell hit in layer $i$ of the \ac{SL}. Of course this is only valid if one assumes simultaneous creation of the ionization in all cells. Since the cosmic muons travel approximately with the speed of light and the chamber measures only $\approx\SI{36}{cm}$ in height (sec.\ \ref{sec:CMS-AachenTestStand-Phase2Setup}), this assumption can be justified.
\FloatBarrier
\paragraph{Pattern finder algorithm}
The algorithm which was implemented in the custom analysis framework to identify the patterns from above should be briefly described here. The algorithm is capable of finding patterns in the recorded cell hits over a full \ac{SL} and is event-loop based. First, all hits are sorted by their timestamp. Then, the algorithm loops through the hit list one by one and stores the current hit in a matrix containing the latest hit of each cell. After the matrix was updated with the current hit, the all possible patterns which contain the most recently added hit are checked, i.e.\ the time difference of the latest cell hits of each possible pattern is calculated. If any of the checked patterns shows a maximum time difference between all hits $\Delta T_\text{max}$ which is not larger than the selected maximum timestamp difference $t_\text{drift,max}$, the pattern is accepted, and all hit information is stored.
\newline Note that because this algorithm checks for new patterns which could be constructed with each new hit, it is possible that multiple patterns are constructed in the event of a noise hit close to the "real" pattern, or in case of high energy particle showers triggering many cells in a region of the chamber. The fact that there may be several patterns very close in time at the same position in the \ac{SL} is revisited in sec.\ \ref{sec:DTData-SLFitGrouping} later.
\FloatBarrier
\paragraph{Applied selection criteria}
This section should present the pattern selection criteria used in this work. Only \ac{SL} patterns with four hits are accepted and not less. In principle (depending on the assumptions) it is possible to also use three-hit patterns for the analysis, however in this thesis this is not the case. Also, regarding the pattern shape, only the patterns with indices $0-5$ are considered. The patterns $6-7$ introduce additional difficulties in the reconstruction without a trigger, and are therefore neglected here. Tab.\ \ref{tab:DTData-PatternSelection} summarizes the pattern selection criteria used in this work.
\begin{table}[htb]
    \centering
    \begin{tabular}{|c|c|c|} 
        \hline
        Parameter & Selection & Details \\ \hline
        Pattern type / shape & Patterns $0-5$ & Fig.\ \ref{fig:DTData-SLPatterns} \\
        Pattern cell hit count $N_\text{hits}$ & $N_\text{hits}=4$ & $-$ \\
        Timestamp difference $\Delta T_\text{max}$ & $\Delta T_\text{max} \leq t_\text{drift,max}\approx\SI{385}{\nano\second}$ & Eq.\ \ref{eq:DTData-PatternTimeWindow} \\ \hline
    \end{tabular}
    \caption{\ac{SL} pattern selection criteria used in this work.}
    \label{tab:DTData-PatternSelection}
\end{table}
\newline The goal of this analysis is to accomplish a track reconstruction algorithm which is simple and can be used for cross-checks with the developed scintillator readout system, and not to conduct precision measurements. Therefore, the reduced selection presented above is considered sufficient. Of course, the algorithm could be extended to both three-hit patterns and different pattern layouts in the future.
\FloatBarrier
\paragraph{Pattern occupancies}
Running the pattern finder algorithm on the recorded \ac{DT} dataset yields a total of \num{81552586} patterns. Fig.\ \ref{fig:DTData-Hist-PatternType} presents the resulting occupancies of the six analyzed pattern types discussed above.
\begin{SCfigure}[50][htb]
    \centering
    \includegraphics[width=0.49\textwidth]{large_combined_data_plots/53_comb__SL_PATTERNS_HIST_pat_type.pdf}
    \caption{Histogram of pattern occupancy binned by pattern type.}
    \label{fig:DTData-Hist-PatternType}
\end{SCfigure}
\newline For all three \acp{SL}, the rates of all pattern types is calculated. Tab.\ \ref{tab:DTData-PatternOccupancies} presents the resulting rates. The observed cumulated pattern rates are around $\approx\SI{700}{\hertz}$ per \ac{SL}. Since the \acp{SL} are approximately similar in area, one would naively expect a similar rate for all of them. Note however due to the distribution of dead cells, some non-uniformity is natural. The lower pattern rate in \ac{SL} 2 might be related to the fact that it has many more dead channels than the other \acp{SL} (tab.\ \ref{tab:DTData-DeadWires}).
\begin{table}[htb]
    \centering  %% VALUES: WITH FULL DATASET
    \begin{tabular}{|c|c|c|c|} 
        \hline
        Pattern & \ac{SL} 1 (phi) & \ac{SL} 2 (theta) & \ac{SL} 3 (phi) \\ \hline
        $0$ & $(213.61 \pm 0.10)\;\si{\hertz}$ & $(181.12 \pm 0.10)\;\si{\hertz}$ & $(204.42 \pm 0.10)\;\si{\hertz}$ \\
        $1$ & $(214.73 \pm 0.10)\;\si{\hertz}$ & $(175.13 \pm 0.09)\;\si{\hertz}$ & $(203.98 \pm 0.10)\;\si{\hertz}$ \\
        $2$ & $(85.67 \pm 0.07)\;\si{\hertz}$ & $(67.23 \pm 0.06)\;\si{\hertz}$ & $(73.74 \pm 0.06)\;\si{\hertz}$ \\
        $3$ & $(81.11 \pm 0.06)\;\si{\hertz}$ & $(66.25 \pm 0.06)\;\si{\hertz}$ & $(73.58 \pm 0.06)\;\si{\hertz}$ \\
        $4$ & $(80.64 \pm 0.06)\;\si{\hertz}$ & $(63.60 \pm 0.06)\;\si{\hertz}$ & $(71.58 \pm 0.06)\;\si{\hertz}$ \\
        $5$ & $(83.39 \pm 0.07)\;\si{\hertz}$ & $(67.23 \pm 0.06)\;\si{\hertz}$ & $(73.58 \pm 0.06)\;\si{\hertz}$ \\ \hline
        Cumulative & $(759.14 \pm 0.20)\;\si{\hertz}$ & $(620.57 \pm 0.18)\;\si{\hertz}$ & $(700.88 \pm 0.19)\;\si{\hertz}$ \\ \hline
    \end{tabular}
    \caption{Observed \ac{SL} pattern rates in the analyzed \ac{DT} dataset.}
    \label{tab:DTData-PatternOccupancies}
\end{table}
\newline It is evident that patterns $(0,1)$, $(2,3)$ and $(4,5)$ exhibit similar rates, which is expected because these patterns are symmetric in their shape (fig.\ \ref{fig:DTData-SLPatterns}). Also, the patterns $(0,1)$ seem to show a significantly higher occupancy as the other pattern types. Since these two patterns generally have a larger angular acceptance than the others, a higher rate is expected. The angular range of the patters is discussed in more detail in sec.\ \ref{sec:DTData-SLPatternFitting-Strategy}, when it comes to track fitting.
\paragraph{Comparison with cosmic rate}
The simulation which was already mentioned above when discussing the cosmic expectation for the \ac{DT} cells is also capable to estimate the cell rate. This is discussed in detail in appendix \ref{sec:Simulation}. For this, the simulated \ac{DT} hits are analyzed with the same pattern finding algorithm (described above) as the chamber data. The expected rates for all analyzed patters are given in tab.\ \ref{tab:Simulation-PatternOccupancies} in the appendix. Details should not be discussed here, but it is evident that the observed pattern rates in data ($\approx\SI{700}{\hertz}$ cumulated per \ac{SL}) are $\approx 30\%$ higher than expected from the simulation ($\approx\SI{540}{\hertz}$ cumulated per \ac{SL}) This phenomenon was already observed similarly for the single cell rates (sec.\ \ref{sec:DTData-CellHits-Occupancy}).
\newline This higher rate could also be explained with the noise hits of the single cells, which can lead to noise-induced patterns, if a coincidence of noise takes place. Also note that a noise hit close to a "real" track pattern leads to additional patters being discovered by the pattern finder algorithm (as already discussed above), since it checks for possible patterns after each event when looping though the drift cell hits. The question which of these patterns are the "real" ones can only be resolved during pattern track fitting and grouping (secs.\ \ref{sec:DTData-SLPatternFitting} and \ref{sec:DTData-SLFitGrouping}). The next section discusses the influence of noise-induced patterns in more detail.

\FloatBarrier
\subsection{Fake patterns}
\label{sec:DTData-SLHitPatterns-FakePatterns}
This section further investigates the influence of noise on the pattern rate by searching for "fake patterns", i.e.\ patterns that can not originate from muon tracks.
\FloatBarrier
\paragraph{Influence of noise hits}
It is clear that the observed noise in the single-cell hits can propagate through the pattern finding algorithm and generate patterns which do not originate from a muon track. While for the track-like patterns these noise-induced patterns mix with the "real" muon track patterns, there are also noise-induced patterns which could not originate from a muon track due to the pattern geometry.\footnote{In principle, there is the possibility that more exotic patterns are formed from real particle crossings. This may be due to showering of high energy particles or due to Coulomb scattering (sec.\ \ref{sec:Interactions-CoulombScattering}) inside the \ac{SL}. These effects are considered to have a small influence, and therefore the fake patterns discussed here are assumed to be generated only by noise.} Analysis of the rate of these fake patterns could therefore provide information on the order of magnitude of the fraction of noise in the track-like patterns.
\FloatBarrier
\paragraph{SL fake patterns}
The selected fake patterns analyzed in this section are presented in fig.\ \ref{fig:DTData-SLFakePatterns}. It is obvious, that one can not create these patters by drawing a straight track through the \ac{SL}.\footnote{In fact, fake patterns $(2,3)$ may originate from a muon track which is diagonal and does not trigger the cells between the two hit clusters. However, this is considered very unlikely and this pattern is also used as a fake pattern here.} The same pattern finding algorithm as before is used on the cell hit data, with the task to now identify these fake patterns.
\begin{figure}[htb]
    \centering
    \includegraphics[width=0.9\textwidth]{SL-FakePatterns.pdf}
    \caption{Analyzed fake patterns which can not originate from a muon track.}
    \label{fig:DTData-SLFakePatterns}
\end{figure}
\FloatBarrier
\paragraph{Fake pattern occupancies}
In total, \num{4055402} fake patterns are found in the recorded dataset. Fig.\ \ref{fig:DTData-Hist-FakePatternType} presents the pattern occupancies for the different fake pattern types.
\begin{SCfigure}[50][htb]
    \centering
    \includegraphics[width=0.49\textwidth]{large_combined_data_plots/53_comb__SL_FAKE_PATTERNS_HIST_pat_type.pdf}
    \caption{Histogram of fake pattern occupancy binned by pattern type.}
    \label{fig:DTData-Hist-FakePatternType}
\end{SCfigure}
\newline Tab.\ \ref{tab:DTData-FakePatternOccupancies} lists the calculated rates for the fake patterns for all \acp{SL}. The observed fake pattern rates per \ac{SL} is $\approx 35 \;\si{\hertz}$, and therefore significantly smaller than the rates of the track-like patterns (which was $\approx\SI{700}{\hertz}$, cf.\ \ref{sec:DTData-SLHitPatterns-TrackPatterns}). This emphasizes that there is actually a track-like signal present in the data and not only noise is measured.
\begin{table}[htb]
    \centering  %% VALUES: WITH FULL DATASET
    \begin{tabular}{|c|c|c|c|}  
        \hline
        Pattern & \ac{SL} 1 (phi) & \ac{SL} 2 (theta) & \ac{SL} 3 (phi) \\ \hline
        $0$ & $(2.57 \pm 0.01)\;\si{\hertz}$ & $(1.56 \pm 0.01)\;\si{\hertz}$ & $(1.74 \pm 0.01)\;\si{\hertz}$ \\
        $1$ & $(2.60 \pm 0.01)\;\si{\hertz}$ & $(1.62 \pm 0.01)\;\si{\hertz}$ & $(1.78 \pm 0.01)\;\si{\hertz}$ \\
        $2$ & $(11.67 \pm 0.02)\;\si{\hertz}$ & $(6.74 \pm 0.02)\;\si{\hertz}$ & $(6.91 \pm 0.02)\;\si{\hertz}$ \\
        $3$ & $(12.15 \pm 0.02)\;\si{\hertz}$ & $(7.90 \pm 0.02)\;\si{\hertz}$ & $(7.67 \pm 0.02)\;\si{\hertz}$ \\
        $4$ & $(8.25 \pm 0.02)\;\si{\hertz}$ & $(4.99 \pm 0.02)\;\si{\hertz}$ & $(6.03 \pm 0.02)\;\si{\hertz}$ \\
        $5$ & $(8.25 \pm 0.02)\;\si{\hertz}$ & $(5.00 \pm 0.02)\;\si{\hertz}$ & $(6.04 \pm 0.02)\;\si{\hertz}$ \\ \hline
        Cumulative & $(45.48 \pm 0.05)\;\si{\hertz}$ & $(27.81 \pm 0.04)\;\si{\hertz}$ & $(30.17 \pm 0.04)\;\si{\hertz}$ \\ \hline
    \end{tabular}
    \caption{Observed \ac{SL} fake pattern rates in the analyzed \ac{DT} dataset.}
    \label{tab:DTData-FakePatternOccupancies}
\end{table}
\newline An approximate symmetry in rate for the symmetric fake patterns $(0,1)$, $(2,3)$ and $(4,5)$ is again visible, similar as for the track-like patterns. However, the symmetric fake pattern groups do not all have the same rate. This may be related to the fact that the different fake patterns include different signal fractions: While it is almost impossible to create the fake patterns $(0,1)$ from a muon track with additional noise, this more likely for the fake patterns $(2,3)$, as they have the same geometry as the track-like patterns $(2,3,4,5)$ with one additional noise hit and one track hit left out. From this argument it becomes more plausible why the fake patterns $(0,1)$ have a higher rate. A slight asymmetry between fake patterns $2$ and $3$ is observed, which is generally not expected. It may be related to the distribution of dead wires in the chamber.
\FloatBarrier
\paragraph{Estimation of noise-induced track-like patterns}
If one assumes the noise to be independent for all drift cells, one can now estimate the fraction of noise in the observed track-like patterns with the results from the fake pattern analysis. With the six fake patterns considered, the noise-induced rate per pattern can be roughly estimated as $\approx\SI{35}{\hertz}/6 \approx 6 \;\si{\hertz}$.
\newline As discussed above, the observed discrepancy between the track-like pattern rates in data and simulation is around $\approx 30\%$ (i.e.\ $\approx\SI{150}{\hertz}$), cumulated for the six analyzed patterns (cf.\ sec.\ \ref{sec:DTData-SLHitPatterns-TrackPatterns}). This is larger than what can be explained with the influence of noise, at least in the presented estimation with the fake patterns.
\newline However, it should be noted that the simulation (appendix \ref{sec:Simulation}) is based significant simplifications and therefore serves more as an order-of-magnitude estimation than a precise tool. To this extent, the comparison between simulation and data is successful, as both track-like pattern rates lie in the same order of magnitude of $\mathcal{O}(10^2\;\si{\hertz})$.





\FloatBarrier
\section{Superlayer pattern fitting}
%\begin{itemize}
%    \item Assume constant drift velocity, straight muon track
%    \item Perform linear fit for all laterality combinations with parameters: Arrival time, position in cell, track inclination
%    \item Select laterality combination with best goodness of fit
%    \item Discuss fit results
%    \item Apply appropriate goodness of fit cut
%\end{itemize}
\label{sec:DTData-SLPatternFitting}
%\section{Superlayer relative time alignment}
%\begin{itemize}
%    \item Mention relative alignment problem of OBDTs
%    \item Show double peak in time difference between SL fits
%    \item Establish relative time correction
%\end{itemize}
When the hits of a found track-like patterns really originate from a passing muon, their timestamps are not independent but correlated according to the muon path through the \ac{SL}. By making assumptions on the muon path and the drift behavior of the ionization electrons in the drift cells, one can make a fit to the timestamps of the track-like patterns. This fit allows the estimation of the muon arrival time, as well as its path (i.e.\ position and direction). Also, by considering the goodness-of-fit variable, it is possible to reject noise-induced track-like patterns, if their timestamps are unlikely to match together to any muon path.
\newline This fitting step is applied on all found track-like \ac{SL} patterns separately, incorporating the timestamps of four cell hits. The combination of the fitted paths to a global muon track is performed later in the analysis (sec.\ \ref{sec:DTData-GlobalTrackBuilding}).

\FloatBarrier
\subsection{Assumptions and fit strategy}
\label{sec:DTData-SLPatternFitting-Strategy}
\FloatBarrier
\paragraph{Approximations and track description}
There are several simplifying assumptions made for the fitting of the hit timestamps of the found patterns. These assumptions are explained in the following, together with the mathematical description of the muon track. Fig.\ \ref{fig:DTData-SLPatternFit-Illustration} illustrates an example muon track with the respective \ac{SL} pattern, and the parameters used in the pattern fit in this section as well as the used local coordinate system.
\begin{figure}[htb]
    \centering
    \hfill
    \begin{subfigure}[c]{0.45\textwidth}
        \centering
        \includegraphics[width=\textwidth]{SL-PatternFit.pdf}
        \caption{Pattern with track and its parameters $x_0$, $\alpha$ and $T_0$}
        \label{fig:DTData-SLPatternFit-Illustration-Track}
    \end{subfigure}
    \hfill
    \begin{subfigure}[c]{0.45\textwidth}
        \centering
        \includegraphics[width=0.7\textwidth]{SL-PatternFit-Cell.pdf}
        \caption{Single cell and laterality $L$ definition}
        \label{fig:DTData-SLPatternFit-Illustration-Cell}
    \end{subfigure}
    \hfill
    \caption{Coordinate system convention and illustration of the parameters of the \ac{SL} pattern fit.}
    \label{fig:DTData-SLPatternFit-Illustration}
\end{figure}
\newline Regarding the muon track, it is assumed to be a straight line, i.e.\ the influence of Coulomb scattering (sec.\ \ref{sec:Interactions-CoulombScattering}) is neglected. In an $x$-$z$-projection (the $x$ axis in this case is perpendicular to the drift cell wire), one can describe the horizontal position of the muon with the following equation:
\begin{equation}
    x_\mu (z) = x_0 + (z-z_{ref}) \cdot \tan\alpha,
    \label{eq:SLPatternFit-LinearTrack}
\end{equation}
where $x_0$ is the reference position of the muon (at height $z=z_{ref}$) and $\alpha$ is the track angle measured from the $z$-axis ($\alpha=0$ corresponds to a vertical track). The local coordinate system is chosen with the origin at the position of the $ly=3$ wire, therefore it is $z_{ref}=z_{\text{wire},ly=3}$. The angle $\alpha$ is measured from the vertical line and is positive towards the left.\footnote{This unusual convention is used here since the $z$ axis points upwards but $(z-z_{ref})$ is negative for the cells below the reference cell.} The parameters $\alpha$ are illustrated in fig.\ \ref{fig:DTData-SLPatternFit-Illustration-Track}.
\newline The drift velocity $v_D$ of the electrons in the cell is assumed to be constant over the whole length. Therefore, the drift distance $x_D$ can simply be related to the drift time $t$ as follows:
\begin{equation}
    x_D = v_D \cdot t.
    \label{eq:SLPatternFit-DriftDistance}
\end{equation}
This is of course an approximation, as the electric field is not completely homogeneous, especially near the wire and not over the full cell height.
\newline To relate the straight muon track to the drift velocities of the cells, one needs to equate the expected $x$-position of the track $x_\mu$ with the wire position $x_{\text{wire},ly}$ plus the drift distance $x_D$. As the wire is at the center of the drift cell and only the (non-negative) drift time $t$ is measured, a left-right ambiguity remains. This ambiguity is described with the laterality factor $L=\pm 1$.\footnote{The sign of the laterality is chosen as follows: $L=-1$ means left ($x<x_\text{wire}$) of the wire, $L=+1$ means right ($x>x_\text{wire}$) of the wire.} For the drift cells of the four layers $ly =0,1,2,3$, one can then write down the following equation:
\begin{equation}
    x_\mu (z=z_{\text{wire},ly}) \overset{!}{=} x_{\text{wire},ly} + L_{ly} \cdot v_D \cdot t_{ly}.
    \label{eq:SLPatternFit-xPosition}
\end{equation}
Here, the vertical position of the drifting electrons is assumed to be the height of the respective wire $z_{\text{wire},ly}$. In reality however, the muon leaves a track of ionization charges over the full height of the cell, not a localized charge at the height of the wire. Especially for larger muon inclinations, this should imply a bias on this fit method. Fig.\ \ref{fig:DTData-SLPatternFit-Illustration-Cell} above illustrates the definition and sign convention used for the laterality.
\newline One further assumes that the ionization events in all cells happen at the same time, which is equivalent to an infinitely fast muon. Given the height of one \ac{SL} of \SI{5.2}{\cm} and the muon travelling at approximately the speed of light, this assumption can be justified. With this assumption one can use a common arrival timestamp $T_0$ of the muon. The recorded timestamps of the cell hits in the four layers $T_{ly}$ can then be related to the drift time $t_{ly}$ for the respective cell as follows:
\begin{equation}
    T_{ly} = T_0 + t_{ly}.
    \label{eq:SLPatternFit-Timestamp}
\end{equation}
\FloatBarrier
\paragraph{Timestamp fit function}
Combining eqs.\ \ref{eq:SLPatternFit-LinearTrack}, \ref{eq:SLPatternFit-xPosition} and \ref{eq:SLPatternFit-Timestamp} yields the following formula for the cell timestamps $T_{ly}$:
\begin{equation}
    T_{ly} = T_0 + \frac{L_{ly}}{v_D} \cdot \left( x_0 + z_{\text{wire},ly} \cdot \tan\alpha - x_{\text{wire},ly} \right).
    \label{eq:SLPatternFit-FitFunction}
\end{equation}
This equation relates the four measured timestamps to the muon track parameters $x_0$ and $\alpha$, as well as the arrival time $T_0$. By fitting the measured timestamps $T_{ly}$ to this function, one can obtain the values for these parameters. Fig.\ \ref{fig:DTData-SLPatternFit-Illustration-Track} from above illustrates the meaning of the three parameters.
\newline Note that in principle one could even use the drift velocity $v_D$ as an additional fit parameter, if one uses patterns with four hits. However, in this thesis the drift velocity is fixed to the known \ac{CMS} reference value of $v_D\approx\SI{54.5}{\micro\meter/\nano\second}$ (eq.\ \ref{eq:DT-DriftVelocity}). In principle, it would be possible to use the value measured with the photo-peak method above (eq.\ \ref{eq:PhotoPeak-DriftVelocityEstimation}), but due to the observed uncertainties of this method, the obtained value (which is lower than the \ac{CMS} value) is not trusted for the reconstruction. However, it is clear that a wrongly assumed drift velocity has consequences for the reconstruction quality.
\FloatBarrier
\paragraph{Fit method and uncertainties}
The fit method is a least-square minimization ($\chi^2$-fit) which uses the layer $ly=0,1,2,3$ as $x$-coordinate and the timestamp $T_{ly}$ as $y$-coordinate. It is clear that when fitting four timestamps to three parameters, the number of degrees of freedom is $N_{df}=1$.
\newline The uncertainty of the timestamps $\sigma_{T,ly}$ for the fit is derived from the digitization uncertainty ($0.22 \;\si{\nano\second}$) and the uncertainty of the applied timing correction obtained from the test pulse run (typically $1-2 \;\si{\nano\second}$, sec.\ \ref{sec:DTData-Testpulse-Calibration}). In addition to these, a generous \SI{5}{\nano\second} uncertainty is assigned to all timestamps which should compensate for the simplified assumption of the drift electrons always being at the height of the wire (see discussion above). With this, the total assigned uncertainty of the cell timestamps is in the order of $6 \;\si{\nano\second}$.
\FloatBarrier
\paragraph{Laterality combinations}
The value of the laterality factors $L_{ly}$ is initially unclear, as there are several allowed laterality configurations of the four cells for each pattern. Tab.\ \ref{tab:DTData-SLPatternFit-AllowedLateralities} in the appendix lists all possibilities for each of the analyzed patterns. As the assumption of one laterality configuration would result in a strong bias for the analysis, the timestamp fit is performed for each combination. Then, only the fit with the best (lowest) goodness-of-fit variable $\chi2/N_{df}$ is kept for the next analysis steps.
\FloatBarrier
\paragraph{Fit parameter bounds and initial values}
The fit parameter values are required to obey to certain bounds, which represent the \ac{DT} chamber properties. The reference position $x_0$ of the track in the \ac{SL} fit is required to be within the reference cell ($ly=3$). Depending on the laterality, it must be within the cell frame and left or right of the wire:
\begin{equation}
    \left.\begin{aligned}
        \SI{-21}{\milli\meter},\;\; & L_{ly=3}=-1 \\
        \SI{0}{\milli\meter},\;\; & L_{ly=3}=+1
    \end{aligned}\right\}
    \leq x_0 \leq
    \left\{\begin{aligned}
        \SI{0}{\milli\meter},\;\; & L_{ly=3}=-1 \\
        \SI{21}{\milli\meter},\;\; & L_{ly=3}=+1
    \end{aligned}\right. .
    \label{eq:SLPatternFit-x0-bounds}
\end{equation}
\newline The track angle $\alpha$ is restricted to the maximum angular acceptance of the selected \ac{SL} patterns (sec.\ \ref{sec:DTData-SLHitPatterns-TrackPatterns}). One can calculate the minimum $ \alpha_{x,min}$ and maximum $\alpha_{x,max}$ allowed track angles for each pattern, which are presented in tab.\ \ref{tab:SLPatternFit-AngularRange} in the appendix. For the track fit, one requires the angle to lie within the limits for the respective pattern type:
\begin{equation}
    \alpha_{min} \leq \alpha \leq \alpha_{max} .
    \label{eq:SLPatternFit-alphax-bounds}
\end{equation}
The arrival time $T_0$ of the track can be directly constrained by the four cell timestamps $T_{ly}$. It is clear that the arrival time must lie before the time of the earliest cell hit and after the time of the latest cell hit subtracted by the maximum possible drift time in the cell $t_{max}$:
\begin{equation}
    \max \left\{ T_{ly} ; \; ly=0,1,2,3 \right\} - t_{max}
    \leq T_0 \leq
    \min \left\{ T_{ly} ; \; ly=0,1,2,3 \right\} .
    \label{eq:SLPatternFit-t0-bounds}
\end{equation}
%In the case of noise-induced patterns it is possible that these bound on $T_0$ are invalid, as the lower bound may become larger than the upper bound. In such a case, the pattern is rejected from the fit.
\newline As starting value for the parameters for the fit, the mean value of the upper and lower bound for each parameter is used.
\FloatBarrier
\paragraph{Cuts on fit results}
In addition to the fit parameter bounds discussed above, there are cuts applied to the fit results. As explained, there is the possibility for noise-induced patterns. By applying a cut on the goodness-of-fit variable $\chi^2/N_{df}$, it is possible to reject these patterns, if their timestamps can hardly be used to construct a straight line track. Additionally, one cuts away those fits which have predicted non-physical drift times. It is clear that the drift time $t_{ly}=T_{ly}-T_0$ in all cells should be smaller than the maximum drift time $t_\text{drift,max}$ which is given by the drift velocity assumption. The cuts applied to the \ac{SL} pattern fits are summarized in tab.\ \ref{tab:DTData-SLPatternFit-Cuts}.
\begin{table}[htb]
    \centering
    \begin{tabular}{|c|c|} 
        \hline
        Variable & Condition \\ \hline
        Goodness-of-fit $\chi^2/N_{df}$ & $\chi^2/N_{df} \leq 10$ \\
        Drift time $t_{ly}$ & $\SI{0}{\nano\second} \leq t_{ly} \leq t_\text{drift,max}\approx\SI{385}{\nano\second}$ \\ \hline
    \end{tabular}
    \caption{Applied cuts to \ac{SL} pattern fit results and the cut flow w.r.t.\ the initial dataset.}
    \label{tab:DTData-SLPatternFit-Cuts}
\end{table}

\FloatBarrier
\subsection{Illustration with example}
\label{sec:DTData-SLPatternFitting-Example}
Fig.\ \ref{fig:DTData-SLPatternFit-EventDisplay-TimestampFit} shows the timestamp fit for one example pattern of the recorded data. Note that the $y$ axis of the plot has very large values, as it directly represents the hit timestamp. The track fit seems to approximate the recorded timestamps reasonably well, with deviations of only some timestamp units, i.e.\ few nanoseconds.
\begin{SCfigure}[50][htb]
    \centering
    \includegraphics[width=0.49\textwidth]{dt_sl_fit/timestamp_fit.pdf}
        \caption{Example of a timestamp $T_{ly}$ fit for one track-like \ac{SL} pattern of type 0. The best laterality combination after trying all combinations is $L_{ly}=[-1,1,-1,1]$ (in the order $ly=0,1,2,3$).}
        \label{fig:DTData-SLPatternFit-EventDisplay-TimestampFit}
\end{SCfigure}
\newline The reconstructed position of the cell hits and the local track in the \ac{SL} for the same example pattern are shown in fig.\ \ref{fig:DTData-SLPatternFit-EventDisplay-CellFit}. The hit position in the cells is calculated from the drift time $t_{ly}=T_{ly}-T_0$ using the fitted arrival time $T_0$ and the selected laterality $L_{ly}$. The fitted track seems to match reasonably well with the estimated hit position in all four cells.
\begin{figure}[htb]
    \centering
    \centering
    \includegraphics[width=0.8\textwidth]{dt_sl_fit/sl_fit.pdf}
    \caption{Reconstructed cell hit positions and fitted track $x_\mu(z)$ for the same example \ac{SL} pattern as fig.\ \ref{fig:DTData-SLPatternFit-EventDisplay-TimestampFit}. Note that the laterality $L_{ly}$ determines the hit position relative to the central wire for the cell in each layer ($L_{ly}=-1$: left, $L_{ly}=+1$: right).}
    \label{fig:DTData-SLPatternFit-EventDisplay-CellFit}
\end{figure}

\FloatBarrier
\subsection{Pattern fit results}
\label{sec:DTData-SLPatternFitting-Results}
\FloatBarrier
\paragraph{Fitted pattern occupancy}
%%% sl patterns: 81552586  --  sl fitted patterns, with cuts: 47959940  --  cut flow: 58.8 %
The track fits are performed for all patterns and the discussed cuts (tab.\ \ref{tab:DTData-SLPatternFit-Cuts}) are applied to the results. With these, the number of patterns after the fit is reduced to \num{47959940} ($\approx 58.8 \%$ of the initial count). The rates of fitted patterns in the different \acp{SL} for all pattern types are presented in tab.\ \ref{tab:DTData-SLPatternFit-PatternOccupancies}.
\begin{table}[htb]
    \centering  %% VALUES: WITH FULL DATASET
    \begin{tabular}{|c|c|c|c|} 
        \hline
        Pattern & \ac{SL} 1 (phi) & \ac{SL} 2 (theta) & \ac{SL} 3 (phi) \\ \hline
        $0$ & $(139.74 \pm 0.08)\;\si{\hertz}$ & $(121.40 \pm 0.08)\;\si{\hertz}$ & $(142.10 \pm 0.09)\;\si{\hertz}$ \\
        $1$ & $(142.32 \pm 0.09)\;\si{\hertz}$ & $(117.61 \pm 0.08)\;\si{\hertz}$ & $(142.54 \pm 0.09)\;\si{\hertz}$ \\
        $2$ & $(41.38 \pm 0.05)\;\si{\hertz}$ & $(31.07 \pm 0.04)\;\si{\hertz}$ & $(37.93 \pm 0.04)\;\si{\hertz}$ \\
        $3$ & $(32.86 \pm 0.04)\;\si{\hertz}$ & $(31.49 \pm 0.04)\;\si{\hertz}$ & $(36.67 \pm 0.04)\;\si{\hertz}$ \\
        $4$ & $(34.99 \pm 0.04)\;\si{\hertz}$ & $(30.17 \pm 0.04)\;\si{\hertz}$ & $(35.81 \pm 0.04)\;\si{\hertz}$ \\
        $5$ & $(36.56 \pm 0.04)\;\si{\hertz}$ & $(31.59 \pm 0.04)\;\si{\hertz}$ & $(37.34 \pm 0.04)\;\si{\hertz}$ \\ \hline
        Cumulative & $(427.84 \pm 0.15)\;\si{\hertz}$ & $(363.32 \pm 0.14)\;\si{\hertz}$ & $(432.40 \pm 0.15)\;\si{\hertz}$ \\ \hline
    \end{tabular}
    \caption{Rates of the track-like \ac{SL} patterns after the track fit, with the described cuts applied.}
    \label{tab:DTData-SLPatternFit-PatternOccupancies}
\end{table}
\newline The rate reduction compared to the set of track-like patterns before the fit and the applied cuts (tab.\ \ref{tab:DTData-PatternOccupancies}) seems to be by a similar factor for all \acp{SL} and pattern types. Note that the fitted pattern rate is reduced to $\approx\SI{400}{\hertz}$ per \ac{SL} through the cuts. The fact that this reduction takes place, emphasizes that there were noise-induced patterns in the dataset which are removed by the selected cut conditions. Of course it is also possible that some real muon patterns are cut with the conditions, as the track fitting is performed with several simplifications and assumptions that were discussed above (sec.\ \ref{sec:DTData-SLPatternFitting-Strategy}).
\newline When comparing the obtained rate to the rate from patterns by simulated atmospheric muon tracks ($\approx\SI{540}{\hertz}$ per \ac{SL}, cf.\ \ref{tab:DTData-PatternOccupancies}), it is clear that the fitted pattern rate is now smaller than the expectation. Here, it must be noted that the simulation does not consider the dead cells present in the real chamber. In any case, after fitting and applying the cuts one can conclude that the obtained fitted pattern rate definitely came closer to the expected rate from simulation.
\FloatBarrier
\paragraph{Goodness-of-fit distribution}
A histogram of the goodness-of-fit parameter $\chi^2/N_{df}$ for the pattern fits after the applied cuts is presented in fig.\ \ref{fig:DTData-Hist-SLPatternFit-chi2}. Generally it can be said that a more steep distribution indicates overall better fit results.
\newline The majority of pattern fits in data after the applied cuts have a good fit quality, with values of $\chi^2/N_{df}<1$. Towards larger goodness-of-fit values (i.e.\ worse fit quality), the number of pattern fits decreases. Given the simple fitting approach used in this work (sec.\ \ref{sec:DTData-SLPatternFitting-Strategy}) is not expected that the pattern fits are always excellent. Also, the contamination of the data with noise hits leads to a less steep distribution.
\begin{SCfigure}[50][htb]
    \centering
    \includegraphics[width=0.49\textwidth]{large_combined_data_plots/53_comb__SL_FITS_AFTERCUTS_HIST_chi2-ndf.pdf}
    \caption{Histogram of goodness-of-fit parameter $\chi^2/N_{df}$ cumulated for all patterns in all \acp{SL}, with the specified cuts applied.}
    \label{fig:DTData-Hist-SLPatternFit-chi2}
\end{SCfigure}
\FloatBarrier
\paragraph{Track parameter distributions}
In the following, histograms of the resulting track parameters $(T_0,\; x_0,\; \tan\;\alpha)$ are discussed. The error bars on the histograms consist of the statistical uncertainty (calculated from the number of bin entries) and the propagated uncertainty from the individual fit parameter values. The error propagation to the histogram is done by shifting the individual values by their error to the left and right, and using the mean change in bin entries compared to the central histogram as uncertainty for each bin. Fig.\ \ref{fig:DTData-Hist-SLPatternfit} presents the histograms for the obtained values of the three fit parameters $T_0$, $x_0$ and $tan_\alpha$ cumulated for all \acp{SL}, after the described cuts were already applied.
\begin{figure}[htb]
    \centering
    \hfill
    \begin{subfigure}[c]{0.49\textwidth}
        \centering
        \includegraphics[width=\textwidth]{large_combined_data_plots/53_comb__SL_FITS_AFTERCUTS_HIST_t0.pdf}
        \caption{Track arrival time $T_0$}
        \label{fig:DTData-Hist-SLPatternfit-t0}
    \end{subfigure}
    \hfill
    \begin{subfigure}[c]{0.49\textwidth}
        \centering
        \includegraphics[width=\textwidth]{large_combined_data_plots/53_comb__SL_FITS_AFTERCUTS_HIST_x0.pdf}
        \caption{Track reference position $x_0$}
        \label{fig:DTData-Hist-SLPatternfit-x0}
    \end{subfigure}
    \hfill
    \begin{subfigure}[c]{0.49\textwidth}
        \centering
        \includegraphics[width=\textwidth]{large_combined_data_plots/53_comb__SL_FITS_AFTERCUTS_HIST_tan_alpha.pdf}
        \caption{Track inclination $\tan\alpha$}
        \label{fig:DTData-Hist-SLPatternfit-tan-alpha}
    \end{subfigure}
    \hfill
    \caption{Histograms of the results for the fit parameters of the pattern fits cumulated for all patterns in all \acp{SL}, with the specified cuts applied.}
    \label{fig:DTData-Hist-SLPatternfit}
\end{figure}
\newline The calculated arrival times $T_0$ (fig.\ \ref{fig:DTData-Hist-SLPatternfit-t0}) seem to be uniformly distributed. This is the expected behavior for atmospheric muons, which should arrive at the chamber at random times.
\newline The reference position of the track $x_0$ (fig.\ \ref{fig:DTData-Hist-SLPatternfit-x0}) shows an approximately uniform distribution over the full cell size of $\pm 21 \;\si{\milli\meter}$, with exception of the region close to the anode wire at $x_0=\SI{0}{\milli\meter}$. As the electric field in this region is hardly homogeneous as assumed in the fit, this deviation is not completely unexpected. Note that the uncertainties of the values close to the wire and at the left and right edges of the drift cell are larger than in the regions in between. A uniform distribution is expected from atmospheric muons, as the impact position of the muon on the chamber is random.
\newline The track inclination $\tan\;\alpha$ (fig.\ \ref{fig:DTData-Hist-SLPatternfit-tan-alpha}) shows an approximately symmetric distribution around $\tan\;\alpha=0$, i.e.\ a vertically impinging track with $\alpha=\SI{0}{\radian}$. The covered angular range is limited by the angular acceptance of the used patterns, which was already discussed in sec.\ \ref{sec:DTData-SLPatternFitting-Strategy}. The shape of the obtained angular distribution is expected from atmospheric muons, as the atmospheric muon tracks should follow a $\propto \cos^2 \theta$ distribution (sec.\ \ref{sec:AtmosphericMuons-Dependencies}). This means that vertical muons the most frequent and the number of muons decreases with increasing inclination. Note that the \ac{SL} pattern fits take place in one 2D slice ($x$-$z$- or $y$-$z$-plane) of space, which means one can not directly translate $\alpha$ to $\theta$. However, the general dependencies of decreasing rate with higher inclination should remain the same.
\newline 
\FloatBarrier
\paragraph{Track parameter uncertainties}
The size of the uncertainties of the obtained fit parameters should be briefly discussed in this paragraph. The uncertainties are extracted from the covariance matrix of the parameters from the individual fits. In fig.\ \ref{fig:DTData-Hist-SLPatternfit-Uncertainties} the appendix, histograms of the uncertainties for all fitted patterns after the applied cuts are presented. Tab.\ \ref{tab:DTData-SLPatternfit-Uncertainties} summarizes the obtained range of uncertainties for the three track parameters.

\FloatBarrier
\paragraph{Drift time distribution}
\FloatBarrier
\paragraph{Comparison with cosmic tracks}





\FloatBarrier
\section{Grouping of patterns within superlayer}
%\begin{itemize}
%    \item Clustering algorithm creates multiple clusters in case of additional nearby hits
%    \item Cope this by grouping together clusters within certain time window within same superlayer
%\end{itemize}
\label{sec:DTData-SLFitGrouping}
\FloatBarrier
\paragraph{SL pattern grouping}
One needs to consider the fact that several patterns close in time may be present in the same \ac{SL}. This possibility arises from the implementation of the pattern finding algorithm, which searches for possible patterns after each hit in the event loop (sec.\ \ref{sec:DTData-SLHitPatterns-TrackPatterns}). An additional cell hit close to a real track pattern can therefore cause an additional pattern. When attempting to build global muon tracks, it is necessary to combine fitted patterns of multiple \acp{SL} (sec.\ \ref{sec:DTData-GlobalTrackBuilding}). Therefore, multiple patterns close in time within the same \ac{SL} will create ambiguities in global track building. The step of pattern grouping takes care of these ambiguities by identifying and handling them.
\newline All patterns of one \ac{SL} are grouped by the arrival timestamp $T_0$ obtained from the pattern fit (sec.\ \ref{sec:DTData-SLPatternFitting}). The time window $\Delta T_\text{0,max}(\text{SL})$ in which the fitted patterns are grouped is chosen as the maximum drift time:
\begin{equation}
    \Delta T_\text{0,max}(\text{SL}) = \max \left\{ \left | T_{0,i,SL} - T_{0,j,SL} \right | ; \; i,j=0,...,N_\text{SL fits}(\text{SL}) \right\} \leq t_\text{drift,max} \approx \SI{385}{\nano\second}.
    \label{eq:DTData-SLFitGrouping-TimeWindow}
\end{equation} 
\FloatBarrier
\paragraph{Grouping results}
\FloatBarrier
\paragraph{Timing difference between SLs}




\FloatBarrier
\section{Global track building}
%\begin{itemize}
%    \item Select cluster groups within time window from all superlayers
%    \item Only use cluster groups with no ambiguity (only one cluster per SL)
%    \item Ensure matching tracks in both phi superlayers
%    \item Calculate 3D track from 2D projections from theta and phi superlayers
%\end{itemize}
\label{sec:DTData-GlobalTrackBuilding}
%\section{Reconstructed muon tracks}
%\begin{itemize}
%    \item Discuss muon rate
%    \item Show angular distributions (with solid angle correction)
%\end{itemize}
%\section{Acceptance correction and angular distributions}
%\begin{itemize}
%    \item Establish geometric simulation of DT chamber with simulated muon tracks
%    \item Derive (rudimentary) geometric acceptance correction from simulation
%    \item Apply acceptance correction to data
%    \item Compare to expected cosmic muon angular distributions
%\end{itemize}
After the local \ac{SL} track fits were performed (sec.\ \ref{sec:DTData-SLPatternFitting}) and the pattern ambiguities in the same \ac{SL} were handled (sec.\ \ref{sec:DTData-SLFitGrouping}), one can finally construct global tracks by combining the information of the different \acp{SL}.

\FloatBarrier
\subsection{Superlayer pattern group matching}
\FloatBarrier
\paragraph{Temporal pattern matching}
The assumption taken during the local \ac{SL} fits, that the muon track creates ionization charges in all cells of the \ac{SL} at the same time (sec.\ \ref{sec:DTData-SLPatternFitting-Strategy}) is extended to all \acp{SL}. Therefore, one would expect the track arrival times for the \ac{SL} fit groups of all three \acp{SL} to be identical.
\newline To account for slight fluctuations, one matches the groups together, if the arrival times $T_{0,SL}$ lie in the following time window $\Delta T_\text{0,max}(\text{global})$:
\begin{equation}
    \begin{aligned}
        \Delta T_\text{0,max}(\text{global}) & = \max \left\{ \left | T_{0,i,SL} - T_{0,j,SL} \right | ; \; i,j=0,...,N_\text{SL fit groups}(\text{SL}) ; \; SL=1,2,3 \right\} \\
        & \leq \frac{t_\text{drift,max}}{2} \approx \SI{193}{\nano\second}.
    \end{aligned}
    \label{eq:DTData-GlobalReco-TimeWindow}
\end{equation} 
\FloatBarrier
\paragraph{Handling of pattern ambiguities}
For cases where a \ac{SL} pattern group contains more than one pattern ($N_\text{pat}>1$), one has an ambiguity in the reconstructed track in this \ac{SL}. In principle one would need to keep all possibilities and may be able to resolve them when combining the \ac{SL} local fit with fits of the other \acp{SL} during global track building. However, this is only possible in the case of the phi \acp{SL}, since only for these one has redundancy and can therefore compare local fits of two different \acp{SL}. Alternatively one could try to resolve the ambiguity by selecting the track fit with the best goodness-of-fit variable.
\newline Since this thesis does only aim at a proof-of-concept measurement of the \ac{DT} chamber and developed scintillator readout system and not at high-precision \ac{DT} reconstruction, this ambiguity should not be considered in detail. Instead, all cases where an ambiguity is present are simply not considered for the global track building, i.e.\ there is a cut applied to the \ac{SL} pattern groups of $N_\text{pat}=1$. This may seem like a strong assumption at first, however the cases with ambiguities are relatively rare. The applied cut only reduces the number of \ac{SL} pattern groups by $\approx 5\%$.
\FloatBarrier
\paragraph{SL selection criteria}
In principle, it would be possible to build a global track when one has at least one \ac{SL} pattern fit in the phi and theta plane. However, in this work one requires a matched pattern in all three \acp{SL}. A global track is therefore only reconstructed when all 12 layers of the chamber are hit and three successful \ac{SL} local pattern fits are available, which have a comparable arrival time.

\FloatBarrier
\subsection{Reconstruction of global track}
\FloatBarrier
\paragraph{Compatibility check of phi projections}
Because of the selection criteria, there are two \ac{SL} phi pattern fits for each global track to be reconstructed. These should first be checked for compatibility and then combined into one track of the \ac{SL} phi projection. The compatibility check has two aspects:
\newline The slope $\tan\alpha_{\phi,i}$ of both local \ac{SL} fits $i=1,2$ is required to not differ more than a chose threshold:
\begin{equation}
    \left| \tan\alpha_{\phi,1},  - \tan\alpha_{\phi,2} \right| \leq 0.1 .
    \label{eq:DTData-GlobalReco-SLPhiMatching-Angle}
\end{equation}
To check whether both fits point to roughly the same position in the chamber, the $x$ coordinates of the intersection points of both local fits $x_{\mu,\phi,i}$ with respect a fixed $z$ position in the chamber are required lie close to each other:
\begin{equation}
    \left| x_{\mu,\phi,1}(z=z_\text{mean}) - x_{\mu,\phi,2}(z=z_\text{mean}) \right| \leq \SI{30}{\milli\meter} .
    \label{eq:DTData-GlobalReco-SLPhiMatching-BasePoint}
\end{equation}
The reference position $z_\text{mean}=\SI{144}{\milli\meter}$ here is approximately the center $z$ position of the \ac{DT} chamber (in the chosen coordinate system, cf.\ sec.\ \ref{sec:CoordinateFrame-TimestampFormat}).
\FloatBarrier
\paragraph{Combination of phi projections}
When the two \ac{SL} phi pattern fits have passed the compatibility checks described above, they can be merged into one track in the phi projection plane. While it in principle would be possible to refit the eight hits of the two phi \acp{SL}, here the parameters are combined in a simpler manner: The slope of the two local track fits $\tan\alpha_{\phi,i}$ is averaged, and the base point of the local track (determined by the $x_{0,\phi,i}$ parameters) is replaced by a new base point which is the average $x$ position of both local tracks at the reference height $z_=z_{mean}$ defined above. Also, the arrival time $T_{0,\phi,i}$ of both tracks is simply averaged.
\FloatBarrier
\paragraph{Global track construction}
The combined phi projection (denoted by index $\phi$) and the theta projection (denoted by index $\phi$) of the track should now be translated to a global track in the 3-dimensional coordinate frame. In the global frame, the track is described with a base point $(x_0,\;y_0,\;z_0)$ and the azimuth $\phi \in [0,\; 2\pi)$ and zenith $\theta \in [0,\;\pi/2)$ angle for a half-sphere.
\newline From the coordinate system convention used in this thesis (sec.\ \ref{sec:CoordinateFrame-TimestampFormat}), it is clear that the \ac{SL} phi projection lies in the $x$-$z$-plane and the \ac{SL} theta projection lies in the $y$-$z$-plane. The transformation for the angles can therefore be calculated as follows:
\begin{equation}
    \phi = \atan2 \left( \frac{\tan \alpha_{\phi}}{\tan \alpha_{\theta}} \right)
\end{equation}
\begin{equation}
    \theta = \arctan \left( \frac{\tan \alpha_{\theta}}{\cos \phi} \right)
\end{equation}
Fig.\ \ref{fig:DTData-GlobalTrack-Illustration} illustrates the relation of the local \ac{SL} projections and the global track variables.
\begin{figure}[htb]
    \centering
    \includegraphics[width=0.45\textwidth]{GlobalTrack.pdf}
    \caption{Illustration of the relation of the \ac{SL} phi projection ($x$-$z$-plane) and \ac{SL} theta ($y$-$z$-plane) local track angles $\alpha_\phi$ and $\alpha_\theta$ to the global track angles $\theta$ and $\phi$ in the used coordinate frame.}
    \label{fig:DTData-GlobalTrack-Illustration}
\end{figure}
\newline The base point of the global track is constructed from the base points of the two projections. The arrival time $T_0$ of the global track is calculated as the average arrival time of the two projections $T_{0,\phi}$ and $T_{0,\theta}$.
\FloatBarrier
\subsection{Global track results}
A
\FloatBarrier
\paragraph{Illustration with event display}
A
\FloatBarrier
\paragraph{Track occupancy}
A
\FloatBarrier
\paragraph{Arrival time}
A
\FloatBarrier
\paragraph{Track position in chamber}
A
\FloatBarrier
\paragraph{Angular distributions}
A
\FloatBarrier
\paragraph{Comparison to cosmics}
A




